\textbf{Phần 1. Trắc nghiệm lựa chọn}
\setcounterpageref{ex}{0}

\begin{ex} %[2D10N51]
Số gồm $3$ trăm, $4$ chục và $5$ đơn vị viết là
\choice
{$354$}
{$543$}
{\True $345$}
{$3405$}
\loigiai{
Số đó là $345$.\\
\textbf{Đáp án của câu hỏi là C}
}
\end{ex}

\begin{ex} %[2D10N51]
Số $528$ đọc là
\choice
{\True Năm trăm hai mươi tám}
{Năm trăm hai mươi}
{Năm trăm hai tám}
{Năm hai tám}
\loigiai{
Đọc là Năm trăm hai mươi tám.\\
\textbf{Đáp án của câu hỏi là A}
}
\end{ex}

\begin{ex} %[2D10H51]
Trong số $764$, chữ số 7 ở hàng nào?
\choice
{Hàng đơn vị}
{\True Hàng trăm}
{Hàng chục}
{Hàng nghìn}
\loigiai{
Chữ số 7 ở hàng trăm.\\
\textbf{Đáp án của câu hỏi là B}
}
\end{ex}

\begin{ex} %[2D10H51]
Trong số $839$, giá trị của chữ số 3 là
\choice
{$3$}
{$300$}
{$39$}
{\True $30$}
\loigiai{
Chữ số 3 ở hàng chục có giá trị là 30.\\
\textbf{Đáp án của câu hỏi là D}
}
\end{ex}

\begin{ex} %[2D10H51]
Số "Bảy trăm linh năm" viết là
\choice
{$750$}
{\True $705$}
{$75$}
{$7005$}
\loigiai{
Viết là $705$.\\
\textbf{Đáp án của câu hỏi là B}
}
\end{ex}

\begin{ex} %[2D10H51]
Số liền sau của số $459$ là
\choice
{$458$}
{$461$}
{$470$}
{\True $460$}
\loigiai{
$459 + 1 = 460$.\\
\textbf{Đáp án của câu hỏi là D}
}
\end{ex}

\begin{ex} %[2D10V51]
Số liền trước của số $600$ là
\choice
{$601$}
{$590$}
{\True $599$}
{$500$}
\loigiai{
$600 - 1 = 599$.\\
\textbf{Đáp án của câu hỏi là C}
}
\end{ex}

\begin{ex} %[2D10V51]
Số bé nhất có ba chữ số khác nhau là
\choice
{$100$}
{\True $102$}
{$120$}
{$123$}
\loigiai{
Số bé nhất có 3 chữ số khác nhau là 102.\\
\textbf{Đáp án của câu hỏi là B}
}
\end{ex}

\begin{ex} %[2D10V51]
Số lớn nhất có ba chữ số khác nhau là
\choice
{\True $987$}
{$999$}
{$989$}
{$897$}
\loigiai{
Số lớn nhất có 3 chữ số khác nhau là 987.\\
\textbf{Đáp án của câu hỏi là A}
}
\end{ex}

\begin{ex} %[2D10V51]
Từ 3 chữ số $1, 2, 3$ lập được bao nhiêu số có ba chữ số khác nhau?
\choice
{$4$ số}
{$8$ số}
{\True $6$ số}
{$3$ số}
\loigiai{
Các số: 123, 132, 213, 231, 312, 321 (có 6 số).\\
\textbf{Đáp án của câu hỏi là C}
}
\end{ex}

\begin{ex} %[2D10V51]
Số gồm 6 trăm và 8 đơn vị viết là
\choice
{$68$}
{$680$}
{$6008$}
{\True $608$}
\loigiai{
Số đó là $608$.\\
\textbf{Đáp án của câu hỏi là D}
}
\end{ex}

\begin{ex} %[2D10C51]
Tìm số có ba chữ số biết chữ số hàng trăm là 4, chữ số hàng chục gấp đôi chữ số hàng trăm và chữ số hàng đơn vị bằng 0.
\choice
{\True $480$}
{$420$}
{$440$}
{$408$}
\loigiai{
Hàng trăm là 4. Hàng chục là $4 \times 2 = 8$. Hàng đơn vị là 0. Số đó là $480$.\\
\textbf{Đáp án của câu hỏi là A}
}
\end{ex}

\textbf{Phần 2. Trắc nghiệm đúng sai}
\setcounterpageref{ex}{0}

\begin{ex} %[2D10H51]
Xét tính đúng sai về cấu tạo số có ba chữ số:
\choiceTF
{\True Số 100 là số bé nhất có ba chữ số.}
{\True Số 999 là số lớn nhất có ba chữ số.}
{Số 305 đọc là Ba trăm năm mươi.}
{\True Số 305 đọc là Ba trăm linh năm.}
\loigiai{
\begin{itemize}
    \item Bé nhất có 3 chữ số là 100: Đúng.
    \item Lớn nhất có 3 chữ số là 999: Đúng.
    \item 305 đọc là Ba trăm linh năm (không phải năm mươi). Ý 3 Sai, Ý 4 Đúng.
\end{itemize}
\begin{itemize}
    \item a) \textbf{Đúng}.
    \item b) \textbf{Đúng}.
    \item c) \textbf{Sai}.
    \item d) \textbf{Đúng}.
\end{itemize}
}
\end{ex}

\begin{ex} %[2D10V51]
Cho số 746. Xét tính đúng sai:
\choiceTF
{\True Chữ số 7 có giá trị là 700.}
{\True Chữ số 4 có giá trị là 40.}
{Chữ số 6 có giá trị là 60.}
{\True Chữ số 6 ở hàng đơn vị nên có giá trị là 6.}
\loigiai{
\begin{itemize}
    \item Chữ số 7 là 700: Đúng.
    \item Chữ số 4 là 40: Đúng.
    \item Chữ số 6 ở hàng đơn vị giá trị 6 (không phải 60). Ý 3 Sai, Ý 4 Đúng.
\end{itemize}
\begin{itemize}
    \item a) \textbf{Đúng}.
    \item b) \textbf{Đúng}.
    \item c) \textbf{Sai}.
    \item d) \textbf{Đúng}.
\end{itemize}
}
\end{ex}

\textbf{Phần 3. Trả lời ngắn}
\setcounterpageref{ex}{0}

\begin{ex} %[2D10H51]
Viết số gồm 5 trăm, 7 chục và 2 đơn vị.
\shortans{572}
\loigiai{
\textbf{Đáp số:} $572$\\
572.
}
\end{ex}

\begin{ex} %[2D10H51]
Viết số gồm 8 trăm và 4 chục.
\shortans{840}
\loigiai{
\textbf{Đáp số:} $840$\\
840.
}
\end{ex}

\begin{ex} %[2D10V51]
Tìm số liền trước của số $350$.
\shortans{349}
\loigiai{
\textbf{Đáp số:} $349$\\
$350 - 1 = 349$.
}
\end{ex}

\begin{ex} %[2D10V51]
Tìm số liền sau của số $789$.
\shortans{790}
\loigiai{
\textbf{Đáp số:} $790$\\
$789 + 1 = 790$.
}
\end{ex}

\begin{ex} %[2D10V51]
Tìm số lớn nhất có ba chữ số khác nhau.
\shortans{987}
\loigiai{
\textbf{Đáp số:} $987$\\
987.
}
\end{ex}

\begin{ex} %[2D10C51]
Tìm tổng các chữ số của số bé nhất có ba chữ số khác nhau ($102$).
\shortans{3}
\loigiai{
\textbf{Đáp số:} $3$\\
$1 + 0 + 2 = 3$.
}
\end{ex}
